\documentclass[../../main.tex]{subfiles}% 注意这里的文件路径不能用 ./main.tex ,否则用latexmk编译子文件会报错
\graphicspath{{\subfix{./image/}}{\subfix{../../image/}}} % 指定图片引用路径,后续可以直接使用图片文件名
% 如果图片存放在上级目录,则使用 ../../image/ 作为备用路径

% 例如:
% \begin{figure}[H]
% \centering
% \includegraphics[width=0.3\textwidth]{图.png}
% \caption{}
% \label{figure:图}
% \end{figure}
% 注意:上述\label{}一定要放在\caption{}之后,否则引用图片序号会只会显示??.

\begin{document}

\section{正则性结果的证明}\label{section:ch07-正则证明}

% SUPPLEMENT: parent=proposition:ch07-7.1.2; kind=proof; 展开四处省略：由\eqref{align:ch07-7.9}令 $l\to\infty$ 得到\eqref{align:ch07-7.157}的控制收敛依据；\eqref{align:ch07-7.158}中公共零测集的可分性论证与 Lebesgue 点；分段 $C^1$ 试验函数 $g_h\zeta$ 的逼近；由\eqref{align:ch07-7.159}取 $\zeta=\rho$ 得测度不等式\eqref{align:ch07-7.15}以及\eqref{align:ch07-7.16}的取极限方式；并在 (ii) 中补出 $f_h$ 的取值、条带 $\Omega_j\times(\tau,\tau+h)$ 上的积分当 $h\to0^+$ 时消失、\eqref{align:ch07-7.161}的计算以及 $j\to\infty$ 时两端的极限
\begin{proof}[\cref{proposition:ch07-7.1.2}的证明]
(i) 只证\eqref{align:ch07-7.15}，\eqref{align:ch07-7.16}的论证与它相同。设 $\widetilde\zeta\in C^1([0,T];C_c^1(\mathbb{R}))$ 满足 $\widetilde\zeta(\cdot,T)=0$ 与 $\widetilde\zeta\geqslant0$。对每个 $l\in[0,\infty)$，\eqref{align:ch07-7.9}（取 $Q_\tau=S$）给出
\begin{align*}
&\iint_S\{[u_r-l]_+\partial_t\widetilde\zeta+H_+(u_r-l)[\phi(u_r)-\phi(l)]\partial_x\widetilde\zeta\}\,\mathrm{d}x\,\mathrm{d}t
+\int_0^T\langle u_s(\cdot,t),\partial_t\widetilde\zeta(\cdot,t)\rangle\,\mathrm{d}t\\
&\qquad\geqslant-\int_\mathbb{R}[u_{0r}-l]_+\widetilde\zeta(x,0)\,\mathrm{d}x-\langle u_{0s},\widetilde\zeta(\cdot,0)\rangle .
\end{align*}
因 $0\leqslant[u_r-l]_+\leqslant u_r\in L^1(S)$、$0\leqslant[u_{0r}-l]_+\leqslant u_{0r}\in L^1(\mathbb{R})$，又 $u_r\geqslant0$ 且 $\phi$ 在 $[0,\infty)$ 上以 $M$ 为 Lipschitz 常数，故在集合 $\{u_r>l\}$ 上有
\begin{align*}
0\leqslant H_+(u_r-l)[\phi(u_r)-\phi(l)]\leqslant|\phi(u_r)-\phi(l)|\leqslant M(u_r-l)\leqslant Mu_r,
\end{align*}
而在 $\{u_r\leqslant l\}$ 上该量等于 $0$；于是对每个 $(x,t)\in S$，当 $l>u_r(x,t)$ 时被积函数为 $0$，由控制收敛定理可在上式中令 $l\to\infty$，其左端前两项与右端第一项均趋于 $0$，得
\begin{align}
\int_0^T\langle u_s(\cdot,t),\partial_t\widetilde\zeta(\cdot,t)\rangle\,\mathrm{d}t\geqslant-\langle u_{0s},\widetilde\zeta(\cdot,0)\rangle . \label{align:ch07-7.157}
\end{align}

对固定的 $\zeta\in C([0,T];C_c(\mathbb{R}))$，函数 $t\mapsto\langle u_s(\cdot,t),\zeta(\cdot,t)\rangle$ 有界可测（界为 $\sup_t\|u_s(\cdot,t)\|_{\mathscr{R}_f(\mathbb{R})}\|\zeta\|_\infty$），故属于 $L^1(0,T)$；由 Lebesgue 微分定理，几乎所有 $\tau\in(0,T)$ 是它的 Lebesgue 点，此时\eqref{align:ch07-7.158}成立。取 $C([0,T];C_c(\mathbb{R}))$ 关于一致范数的可数稠密子集 $\{\zeta_j\}$，把各 $\zeta_j$ 对应的例外零测集合并为 $N$；对一般 $\zeta$，用 $\|\zeta_j-\zeta\|_\infty$ 控制\eqref{align:ch07-7.158}两端之差（两端关于 $\zeta$ 均以 $\sup_t\|u_s(\cdot,t)\|_{\mathscr{R}_f(\mathbb{R})}\|\zeta\|_\infty$ 为界），即得：存在零测集 $N\subseteq(0,T)$，使对一切 $\zeta\in C([0,T];C_c(\mathbb{R}))$ 与 $\tau\in(0,T)\setminus N$，
\begin{align}
\lim_{h\to0^+}\frac1h\int_\tau^{\tau+h}\langle u_s(\cdot,t),\zeta(\cdot,t)\rangle\,\mathrm{d}t=\langle u_s(\cdot,\tau),\zeta(\cdot,\tau)\rangle . \label{align:ch07-7.158}
\end{align}

设 $t_1,t_2\in(0,T)\setminus N$，$0<t_1<t_2<T$，并设 $\zeta\in C^1([0,T];C_c^1(\mathbb{R}))$。对 $h\in(0,\min\{t_2-t_1,T-t_2\})$ 置
\begin{align*}
g_h(t)\triangleq\frac1h(t-t_1)\chi_{[t_1,t_1+h]}(t)+\chi_{(t_1+h,t_2)}(t)+\frac1h(t_2+h-t)\chi_{[t_2,t_2+h]}(t).
\end{align*}
则 $0\leqslant g_h\leqslant1$；$g_h=0$ 在 $[0,t_1]\cup[t_2+h,T]$ 上，$g_h=1$ 在 $[t_1+h,t_2]$ 上；$g_h$ 在 $(t_1,t_1+h)$ 与 $(t_2,t_2+h)$ 外导数几乎处处为 $0$，在这两个区间上分别等于 $1/h$ 与 $-1/h$。取 $C^1$ 函数列 $\{g_h^k\}$ 使 $0\leqslant g_h^k\leqslant1$、$g_h^k\to g_h$ 在 $[0,T]$ 上一致、$(g_h^k)'\to g_h'$ 在 $L^1(0,T)$ 中，把\eqref{align:ch07-7.157}用于试验函数 $g_h^k\zeta\geqslant0$ 并令 $k\to\infty$（各项关于 $\zeta$ 与 $\partial_t\zeta$ 的对偶配对绝对收敛），得\eqref{align:ch07-7.157}对 $g_h\zeta$ 成立。于是由 $\partial_t(g_h\zeta)=g_h'\zeta+g_h\partial_t\zeta$ 与 $g_h(0)=0$，
\begin{align*}
&\frac1h\int_{t_1}^{t_1+h}\langle u_s(\cdot,t),\zeta(\cdot,t)\rangle\,\mathrm{d}t-\frac1h\int_{t_2}^{t_2+h}\langle u_s(\cdot,t),\zeta(\cdot,t)\rangle\,\mathrm{d}t\\
&\qquad+\int_{t_1}^{t_2}g_h(t)\langle u_s(\cdot,t),\partial_t\zeta(\cdot,t)\rangle\,\mathrm{d}t\geqslant0 .
\end{align*}
令 $h\to0^+$：前两项由\eqref{align:ch07-7.158}分别收敛于 $\langle u_s(\cdot,t_1),\zeta(\cdot,t_1)\rangle$ 与 $\langle u_s(\cdot,t_2),\zeta(\cdot,t_2)\rangle$，第三项由 $g_h\to\chi_{[t_1,t_2]}$ 与控制收敛定理收敛于 $\int_{t_1}^{t_2}\langle u_s(\cdot,t),\partial_t\zeta(\cdot,t)\rangle\,\mathrm{d}t$。整理得
\begin{align}
\langle u_s(\cdot,t_2),\zeta(\cdot,t_2)\rangle\leqslant\int_{t_1}^{t_2}\langle u_s(\cdot,t),\partial_t\zeta(\cdot,t)\rangle\,\mathrm{d}t+\langle u_s(\cdot,t_1),\zeta(\cdot,t_1)\rangle . \label{align:ch07-7.159}
\end{align}

在\eqref{align:ch07-7.159}中取 $\zeta(x,t)=\rho(x)$，其中 $\rho\in C_c^1(\mathbb{R})$，$\rho\geqslant0$。则 $\partial_t\zeta=0$，故对一切 $t_1,t_2\in(0,T)\setminus N$，$t_1<t_2$，有 $\langle u_s(\cdot,t_2),\rho\rangle\leqslant\langle u_s(\cdot,t_1),\rho\rangle$；由 $C_c^1(\mathbb{R})$ 中的非负函数在 $C_c(\mathbb{R})$ 的非负元中稠密，该不等式对一切 $\rho\in C_c(\mathbb{R})$、$\rho\geqslant0$ 成立，这正是\eqref{align:ch07-7.15}。至于\eqref{align:ch07-7.16}：把\eqref{align:ch07-7.159}中的 $t_1$ 换成 $0$，即用初值项 $-\langle u_{0s},\zeta(\cdot,0)\rangle$ 代替 $\langle u_s(\cdot,t_1),\zeta(\cdot,t_1)\rangle$（此时 $g_h$ 的左端点平移至 $0$，$g_h(0)=1$），再取 $\zeta=\rho$，得 $\langle u_s(\cdot,t),\rho\rangle\leqslant\langle u_{0s},\rho\rangle$ 对几乎所有 $t\in(0,T)$ 与一切非负 $\rho\in C_c(\mathbb{R})$，即\eqref{align:ch07-7.16}。

(ii) 如\eqref{align:ch07-7.158}那样论证（把 $u_s$ 换成 $u$），存在零测集 $N\subseteq(0,T)$，使对一切 $\rho\in C_c(\mathbb{R})$ 与 $\tau\in(0,T)\setminus N$，
\begin{align}
\lim_{h\to0^+}\frac1h\int_\tau^{\tau+h}\langle u(\cdot,t),\rho\rangle\,\mathrm{d}t=\langle u(\cdot,\tau),\rho\rangle . \label{align:ch07-7.160}
\end{align}
置 $\Omega_j\triangleq[-j-1,-j]\cup[j,j+1]$，$Q_{j,\tau}\triangleq\Omega_j\times(0,\tau)$（$\tau\in(0,T)$，$j\in\mathbb{N}$）。设 $\{\rho_j\}\subseteq C_c^1(\mathbb{R})$ 满足 $\rho_j=1$ 在 $[-j,j]$、$\operatorname{supp}\rho_j\subseteq[-j-1,j+1]$、$0\leqslant\rho_j\leqslant1$ 与 $|\rho_j'|\leqslant2$ 在 $\mathbb{R}$ 上。设 $0<h<\tau$ 并置
\begin{align*}
f_h(t)\triangleq\chi_{[0,\tau)}(t)+\frac1h(\tau+h-t)\chi_{[\tau,\tau+h]}(t),
\end{align*}
故 $f_h=1$ 在 $[0,\tau]$ 上，$f_h$ 在 $[\tau,\tau+h]$ 上线性下降至 $0$，且 $f_h'=-(1/h)\chi_{(\tau,\tau+h)}$ 几乎处处。在\eqref{align:ch07-7.7}中取 $\zeta(x,t)=f_h(t)\rho_j(x)$（取 $Q_\tau=S$；$f_h$ 分段 $C^1$，以 $C^1$ 函数逼近后取极限即得该试验函数的合法性，且 $f_h(T)=0$）。因 $\partial_t\zeta=f_h'\rho_j$、$\partial_x\zeta=f_h\rho_j'$，\eqref{align:ch07-7.7}给出
\begin{align*}
\iint_{Q_{j,\tau}}f_h(t)\phi(u_r)(x,t)\rho_j'(x)\,\mathrm{d}x\,\mathrm{d}t=\frac1h\int_\tau^{\tau+h}\langle u(\cdot,t),\rho_j\rangle\,\mathrm{d}t-\langle u_0,\rho_j\rangle ;
\end{align*}
严格地说，\eqref{align:ch07-7.7}中的 $x$ 积分区域为 $\Omega_j\times(0,T)$，但 $f_h\equiv1$ 在 $(0,\tau)$ 上，且条带 $\Omega_j\times(\tau,\tau+h)$ 上的积分当 $h\to0^+$ 时趋于 $0$，故取极限后的恒等式不变。令 $h\to0^+$ 并利用\eqref{align:ch07-7.160}（右端第一项收敛于 $\langle u(\cdot,\tau),\rho_j\rangle$）与 $f_h\to\chi_{[0,\tau)}$ 的控制收敛定理，得
\begin{align}
\iint_{Q_{j,\tau}}\phi(u_r)(x,t)\rho_j'(x)\,\mathrm{d}x\,\mathrm{d}t=\langle u(\cdot,\tau),\rho_j\rangle-\langle u_0,\rho_j\rangle . \label{align:ch07-7.161}
\end{align}
因 $|\phi(u_r)|\leqslant Mu_r$（见\eqref{align:ch07-7.1}）且 $u_r\in L^\infty(0,T;L^1(\mathbb{R}))$，由 $\operatorname{supp}\rho_j'\subseteq\Omega_j$ 与 $|\rho_j'|\leqslant2$ 得
\begin{align*}
\left|\iint_{Q_{j,\tau}}\phi(u_r)\rho_j'\,\mathrm{d}x\,\mathrm{d}t\right|\leqslant2\int_0^\tau\int_{\{|x|\geqslant j\}}|\phi(u_r)|\,\mathrm{d}x\,\mathrm{d}t\to0\quad(j\to\infty).
\end{align*}
另一方面，$\rho_j\to1$ 在 $\mathbb{R}$ 上处处成立且 $0\leqslant\rho_j\leqslant1$，故由控制收敛定理，\eqref{align:ch07-7.161}右端收敛于 $\langle u(\cdot,\tau),1\rangle-\langle u_0,1\rangle=\|u(\cdot,\tau)\|_{\mathscr{R}_f(\mathbb{R})}-\|u_0\|_{\mathscr{R}_f(\mathbb{R})}$。在\eqref{align:ch07-7.161}中令 $j\to\infty$，即得\eqref{align:ch07-7.17}。结论得证。
\end{proof}

% SUPPLEMENT: parent=proposition:ch07-7.2.1; kind=proof; 展开四处：$Y$ 的三变量记号与“绝对连续部分的密度是函数故无原子”的用法；由\eqref{align:ch07-7.163}的分布导数属于 $\mathscr{R}_f(\mathbb{R})$ 与 $Y(\cdot,t)\in L^1$ 得 $Y(\cdot,t)\in BV(\mathbb{R})$ 及全变差估计；(7.164) 中 $\rho_j'$ 的取值与 $j\to\infty$ 时两端极限的依据；由 (7.162) 推出 (7.19) 时 $u_s(\cdot,t)(\{x_0\})>0$ 与等待时间定义\eqref{align:ch07-7.18}的关系
\begin{proof}[\cref{proposition:ch07-7.2.1}的证明]
对 $0\leqslant t_1\leqslant t_2\leqslant T$ 与 $x\in\mathbb{R}$ 置
\begin{align*}
Y(x,t_1,t_2)\triangleq\int_{t_1}^{t_2}\phi(u_r(x,s))\,\mathrm{d}s,\qquad Y(x,t)\triangleq Y(x,0,t).
\end{align*}
证明下列断言。

\paragraph{结论 1.}
对几乎所有 $t\in(0,T)$，映射 $x\mapsto Y(x,t)$ 属于 $BV(\mathbb{R})$；且对任意 $a,b\in\mathbb{R}$，$a\leqslant b$，有
\begin{align}
u(\cdot,t)([a,b])-u_0([a,b])=Y(a^-,0,t)-Y(b^+,0,t). \label{align:ch07-7.162}
\end{align}

取 $a=b=x_0$。因 $u(\cdot,t)$ 与 $u_0$ 的绝对连续部分的密度是函数，其单点质量为零，故 $u(\cdot,t)(\{x_0\})=u_s(\cdot,t)(\{x_0\})$ 与 $u_0(\{x_0\})=u_{0s}(\{x_0\})$。于是由\eqref{align:ch07-7.162}，对几乎所有 $0\leqslant t\leqslant T$，
\begin{align*}
u_s(\cdot,t)(\{x_0\})=u_{0s}(\{x_0\})+Y(x_0^-,0,t)-Y(x_0^+,0,t)\geqslant u_{0s}(\{x_0\})-2\|\phi\|_\infty t,
\end{align*}
其中用到 $0\leqslant\phi(u_r)\leqslant\|\phi\|_\infty$（$\phi$ 有界且 $u_r\geqslant0$），故 $|Y(x_0^\pm,0,t)|\leqslant\|\phi\|_\infty t$。于是当 $t\in(0,u_{0s}(\{x_0\})/(2\|\phi\|_\infty))$ 时 $u_s(\cdot,t)(\{x_0\})>0$，特别地 $u_s(\cdot,t)\neq0$；按\eqref{align:ch07-7.18}，这样的 $t$ 不属于使“$u_s(\cdot,t)=0$ 且 $u_r(\cdot,t)\in L^\infty(\mathbb{R})$ 对几乎所有 $t\in(\tau,T)$”成立的那些 $\tau$ 的集合，故等待时间 $t_0$ 不小于 $u_{0s}(\{x_0\})/(2\|\phi\|_\infty)$。由此得\eqref{align:ch07-7.19}。

余下证结论 1。如\eqref{align:ch07-7.161}那样论证（把 $\rho_j$ 换成 $\rho$，并保留 $(0,\tau)$ 上的时间积分），得：对几乎所有 $\tau\in(0,T)$ 与一切 $\rho\in C_c^1(\mathbb{R})$，
\begin{align}
\int_\mathbb{R}Y(x,0,\tau)\rho'(x)\,\mathrm{d}x=\langle u(\cdot,\tau),\rho\rangle-\langle u_0,\rho\rangle . \label{align:ch07-7.163}
\end{align}
故分布导数 $\partial_xY(\cdot,0,\tau)$ 属于 $\mathscr{R}_f(\mathbb{R})$，且
\begin{align*}
\|\partial_xY(\cdot,0,\tau)\|_{\mathscr{R}_f(\mathbb{R})}\leqslant\|u(\cdot,\tau)\|_{\mathscr{R}_f(\mathbb{R})}+\|u_0\|_{\mathscr{R}_f(\mathbb{R})}.
\end{align*}
另一方面，由\eqref{align:ch07-7.1}，$|Y(\cdot,0,t)|\leqslant M\int_0^tu_r(\cdot,s)\,\mathrm{d}s\in L^1(\mathbb{R})$，故对几乎所有 $t\in(0,T)$，$Y(\cdot,t)\in L^1(\mathbb{R})$。$L^1(\mathbb{R})$ 中分布导数属于 $\mathscr{R}_f(\mathbb{R})$ 的函数是有界变差函数，其全变差不超过该测度的全变差，故对几乎所有 $t\in(0,T)$，$Y(\cdot,t)\in BV(\mathbb{R})$。

为证\eqref{align:ch07-7.162}，由截断函数的逼近（\eqref{align:ch07-7.163}两端关于 $\rho$ 连续，左端只依赖 $\rho'$），可在\eqref{align:ch07-7.163}中取 $\tau=t$ 与 $\rho=\rho_j$，
\begin{align*}
\rho_j(x)\triangleq j\left(x-a+\frac1j\right)\chi_{[a-1/j,a]}(x)+\chi_{(a,b)}(x)+j\left(b+\frac1j-x\right)\chi_{[b,b+1/j]}(x),
\end{align*}
其中 $a,b\in\mathbb{R}$，$a<b$，$j\in\mathbb{N}$。则 $\rho_j'=j$ 在 $(a-1/j,a)$ 上、$\rho_j'=0$ 在 $(a,b)$ 上、$\rho_j'=-j$ 在 $(b,b+1/j)$ 上，于是
\begin{align}
\langle u(\cdot,\tau),\rho_j\rangle-\langle u_0,\rho_j\rangle=j\int_{a-1/j}^aY(x,0,\tau)\,\mathrm{d}x-j\int_b^{b+1/j}Y(x,0,\tau)\,\mathrm{d}x . \label{align:ch07-7.164}
\end{align}
在\eqref{align:ch07-7.164}中令 $j\to\infty$。因 $0\leqslant\rho_j\leqslant1$ 且 $\rho_j\to\chi_{[a,b]}$ 在 $\mathbb{R}$ 上处处成立，由控制收敛定理，左端趋于 $u(\cdot,\tau)([a,b])-u_0([a,b])$；因 $Y(\cdot,0,\tau)$ 为 $BV$ 函数，其在 $a$ 处的左极限与在 $b$ 处的右极限存在，且收缩区间上的平均值收敛于该单侧极限，故右端趋于 $Y(a^-,0,\tau)-Y(b^+,0,\tau)$。即得\eqref{align:ch07-7.162}。结论 1 与\eqref{align:ch07-7.19}得证。
\end{proof}

\begin{remark}\phantomsection\label{remark:ch07-7.6.1}
\cref{proposition:ch07-7.2.1}的证明事实上依赖于一维双曲守恒律的熵解与 Hamilton--Jacobi 方程的粘性解之间的对应（关于这一主题见文献~\cite{ref21,ref20}）。

该对应在此处的确切形式如下。置 $v(x,t)\triangleq\int_{x_0}^xu(y,t)\,\mathrm{d}y$（测度意义下的积分，$x_0\in\mathbb{R}$）。由\eqref{align:ch07-7.163}，$\partial_xY(\cdot,0,t)=u(\cdot,t)-u_0$，故 $Y(x,0,t)=v(x,t)-v(x,0)$，即 $Y$ 是粘性解 $v$ 的时间增量；而 $v$ 满足 Hamilton--Jacobi 方程 $\partial_tv+\phi(\partial_xv)=0$，其粘性解不等式与 $u=\partial_xv$ 的熵不等式\eqref{align:ch07-7.8}互相等价。\eqref{align:ch07-7.162}正是这一对应的定量形式：右端 $Y(a^-,0,t)-Y(b^+,0,t)$ 是 $v(\cdot,t)-v(\cdot,0)$ 在区间 $[a,b]$ 两端的一侧跳跃，左端是 $u(\cdot,t)$ 在该区间上的质量与初值质量之差；取 $a=b=x_0$ 时该式把 $u_s(\cdot,t)$ 在单点的质量表示为 $Y$ 的跳跃，这正是等待时间估计\eqref{align:ch07-7.19}的来源。
\end{remark}

% SUPPLEMENT: parent=theorem:ch07-7.3.3; kind=proof; 补出“只需在 $\phi''<0$、$\phi'>0$、$\phi(0)=0$ 的假设下证明”的理由：由\cref{remark:ch07-7.1.1}的二分法，情形 (b) 经 $\psi\triangleq-\phi$ 与反射 $x\mapsto-x$ 化归情形 (a)，并逐项核验 (A2')、熵通量随反射的变号与两条结论在该反射下的不变性；同时说明 (V$n^\varepsilon$) 与 (P$n^\varepsilon$) 的差别以及可继承的估计
只需在 (A2') 成立且 $\phi''<0$、$\phi'>0$ 在 $(0,\infty)$、$\phi(0)=0$ 的假设下证明\cref{theorem:ch07-7.3.3}与\cref{proposition:ch07-7.3.5}。由\cref{remark:ch07-7.1.1}，(A2') 蕴含两种互斥情形：情形 (a) $\alpha\phi+\beta>0$ 且 $\phi''<0$ 在 $[0,\infty)$；情形 (b) $\alpha\phi+\beta<0$ 且 $\phi''>0$ 在 $[0,\infty)$。在情形 (b) 中置 $\psi\triangleq-\phi$ 并作反射 $x\mapsto-x$。则 $\psi(0)=0$、$\psi''=-\phi''<0$ 与 $\psi'>0$ 在 $(0,\infty)$，且
\begin{align*}
\psi''(y)[\alpha\psi(y)-\beta]=(-\phi''(y))[-(\alpha\phi(y)+\beta)]=\phi''(y)[\alpha\phi(y)+\beta]\leqslant-[\phi'(y)]^2=-[\psi'(y)]^2<0,
\end{align*}
故 $\psi$ 满足 (A2') 的情形 (a)（常数取 $\alpha$ 与 $-\beta$）。此外，反射 $x\mapsto-x$ 把 $\partial_tu+\partial_x[\phi(u)]=0$ 变为 $\partial_tu+\partial_x[\psi(u)]=0$，并把熵通量 $F_l(y)=\operatorname{sgn}(y-l)[\phi(y)-\phi(l)]$ 变为 $\operatorname{sgn}(y-l)[\psi(y)-\psi(l)]=-F_l(y)$，后者与 $\partial_x$ 的变号相消，故 $\psi$-流对应的熵解正是 $\phi$-流对应的熵解的反射；又\cref{theorem:ch07-7.3.3}的两条结论（在一侧点处的相容条件与 $u=u_r$ 的局部有界性）与\cref{proposition:ch07-7.3.5}的结论（奇异部分支撑集为零测集）在该反射下不变。故只需考虑情形 (a)。此时 $\alpha\phi+\beta>0$ 在 $[0,\infty)$（因 $\phi''[\alpha\phi+\beta]<0$ 与 $\phi''<0$），由 $\phi(0)=0$ 得 $\beta>0$。

考虑 (Pn) 的下列正则化问题（与\cref{section:ch07-逼近问题}中所用的 (P$n^\varepsilon$) 不同）：
\begin{align*}
\begin{cases}
\partial_t y_n^\varepsilon+\partial_x[\phi(y_n^\varepsilon)]=\varepsilon\,\partial_{xx}[\phi(y_n^\varepsilon)] & \text{在 }S,\\
y_n^\varepsilon=u_{0n}^\varepsilon & \text{在 }\mathbb{R}\times\{0\},
\end{cases}
\qquad(Vn^\varepsilon)
\end{align*}
其中 $\{u_{0n}^\varepsilon\}$ 满足\eqref{align:ch07-7.46}--\eqref{align:ch07-7.47}。\cref{section:ch07-逼近问题}中为 (P$n^\varepsilon$) 回顾的存在、唯一与正则性结果，以及\cref{lemma:ch07-7.4.1}的先验估计与\cref{lemma:ch07-7.4.5}的收敛结果，对 (V$n^\varepsilon$) 的解同样成立（见文献~\cite{ref63}）；其中\eqref{align:ch07-7.49}--\eqref{align:ch07-7.52}只用到 $\phi$ 在 $y_n^\varepsilon$ 的值域上的 Lipschitz 性与初值界，\eqref{align:ch07-7.66}--\eqref{align:ch07-7.69}只用到 $\{y_n^\varepsilon\}$ 在 $L^\infty(S)\cap L^\infty(0,T;L^1(\mathbb{R}))$ 中的一致有界性与 $x$ 方向的一致连续性（参见\eqref{align:ch07-7.52}）。族 $\{y_n^\varepsilon\}$ 的若干性质是下列两个引理的内容。

\begin{lemma}\label{lemma:ch07-7.6.1}
设 $\phi$ 满足 (A2')，即 $\phi\in C^\infty([0,\infty))$、$\lim_{y\to\infty}\phi(y)/y=0$，且存在常数 $\alpha\geqslant-1$ 与 $\beta\in\mathbb{R}$，使
\begin{align*}
\phi''(y)[\alpha\phi(y)+\beta]\leqslant-[\phi'(y)]^2<0\quad\text{对一切 }y\in[0,\infty);
\end{align*}
并设 $\phi(0)=0$，$\phi''<0$ 与 $\phi'>0$ 在 $(0,\infty)$。设 $n\in\mathbb{N}$，$T>0$，$S=\mathbb{R}\times(0,T)$；设 $u_{0n}\in L^1(\mathbb{R})\cap L^\infty(\mathbb{R})$ 非负且满足\eqref{align:ch07-7.44}--\eqref{align:ch07-7.45}，$\{u_{0n}^\varepsilon\}\subseteq C_c^\infty(\mathbb{R})$ 非负且满足\eqref{align:ch07-7.46}--\eqref{align:ch07-7.47}；设 $u_n$ 为\cref{proposition:ch07-7.4.7}给出的问题 (Pn) 的唯一熵解。则存在数列 $\{\varepsilon_j\}\subseteq(0,\infty)$，$\varepsilon_j\to0^+$，以及 (V$n^\varepsilon$) 在 $\varepsilon=\varepsilon_j$ 时的解 $y_n^{\varepsilon_j}\in L^\infty(S)\cap L^\infty(0,T;L^1(\mathbb{R}))$，使
\begin{align}
y_n^{\varepsilon_j}\stackrel{*}{\rightharpoonup}u_n\ \text{在 }L^\infty(S),\qquad
y_n^{\varepsilon_j}\to u_n\ \text{在 }L^1(R_L)\text{ 对一切 }L>0. \label{align:ch07-7.165}
\end{align}
\end{lemma}

% SUPPLEMENT: parent=lemma:ch07-7.6.1; kind=proof; 补出三处：(V$n^\varepsilon$) 与 (P$n^\varepsilon$) 共享\eqref{align:ch07-7.66}--\eqref{align:ch07-7.69}的理由；以 $E'(y)\zeta$ 乘 (V$n^\varepsilon$) 所得熵不等式的分部积分计算（含 $\varepsilon_j\iint F'(y)\partial_x y\,\partial_x\zeta$ 的形式与它的 $O(\sqrt{\varepsilon_j})$ 估计）；以及由凸函数 $E$ 的熵不等式过渡到 $E_l(y)=|y-l|$ 的光滑凸逼近
\begin{proof}
由\eqref{align:ch07-7.66}与\eqref{align:ch07-7.69}（二者对 (V$n^\varepsilon$) 的解同样成立），存在子列 $\{y_n^{\varepsilon_j}\}$（$\varepsilon_j\to0^+$）与函数 $y_n\in L^\infty(S)\cap L^\infty(0,T;L^1(\mathbb{R}))$，使\eqref{align:ch07-7.165}成立。此外，设 $E\in C^2([0,\infty))$ 为凸函数，$F\in C^2([0,\infty))$ 满足 $F'=E'\phi'$，$\zeta$ 如\cref{definition:ch07-Pn熵解}（即 $\zeta\in C^1([0,T];C_c^1(\mathbb{R}))$，$\zeta(\cdot,T)=0$，$\zeta\geqslant0$）。以 $E'(y_n^{\varepsilon_j})\zeta$ 乘 (V$n^\varepsilon$)（取 $\varepsilon=\varepsilon_j$）并在 $S$ 上分部积分：由
\begin{align*}
\partial_t[E(y)]+\partial_x[F(y)]=E'(y)\{\partial_ty+\partial_x[\phi(y)]\},\qquad
E'(y)\partial_{xx}[\phi(y)]=\partial_x[F'(y)\partial_xy]-E''(y)\phi'(y)(\partial_xy)^2,
\end{align*}
得
\begin{align*}
&\iint_S\{E(y_n^{\varepsilon_j})\partial_t\zeta+F(y_n^{\varepsilon_j})\partial_x\zeta\}\,\mathrm{d}x\,\mathrm{d}t+\int_\mathbb{R}E(u_{0n}^{\varepsilon_j})\zeta(x,0)\,\mathrm{d}x\\
&\qquad=\varepsilon_j\iint_SF'(y_n^{\varepsilon_j})\,\partial_x y_n^{\varepsilon_j}\,\partial_x\zeta\,\mathrm{d}x\,\mathrm{d}t
+\varepsilon_j\iint_SE''(y_n^{\varepsilon_j})\phi'(y_n^{\varepsilon_j})(\partial_x y_n^{\varepsilon_j})^2\zeta\,\mathrm{d}x\,\mathrm{d}t\\
&\qquad\geqslant\varepsilon_j\iint_SF'(y_n^{\varepsilon_j})\,\partial_x y_n^{\varepsilon_j}\,\partial_x\zeta\,\mathrm{d}x\,\mathrm{d}t,
\end{align*}
最后一个不等号用到 $E''\geqslant0$ 与 $\phi'>0$。右端的 $\varepsilon_j$ 项当 $j\to\infty$ 时趋于 $0$：由 Cauchy--Schwarz 不等式与熵耗散估计 $\varepsilon_j\iint_SE''\phi'(\partial_x y_n^{\varepsilon_j})^2\zeta\leqslant C$（其证明与\eqref{align:ch07-7.59}相同），得
\begin{align*}
\left|\varepsilon_j\iint_SF'(y_n^{\varepsilon_j})\,\partial_x y_n^{\varepsilon_j}\,\partial_x\zeta\,\mathrm{d}x\,\mathrm{d}t\right|
\leqslant C'\sqrt{\varepsilon_j}\left(\varepsilon_j\iint_SE''\phi'(\partial_x y_n^{\varepsilon_j})^2\zeta\,\mathrm{d}x\,\mathrm{d}t\right)^{1/2}\to0 .
\end{align*}
令 $j\to\infty$，并如\cref{proposition:ch07-7.4.7}的证明那样论证（用\eqref{align:ch07-7.69}得 $y_n^{\varepsilon_j}\to y_n$ 在 $S$ 中几乎处处、用\eqref{align:ch07-7.47}得初值项的收敛、用 $\{y_n^{\varepsilon_j}\}$ 的 $L^\infty$ 界与控制收敛定理处理非线性项），得
\begin{align*}
\iint_S[E(y_n)\partial_t\zeta+F(y_n)\partial_x\zeta]\,\mathrm{d}x\,\mathrm{d}t\geqslant-\int_\mathbb{R}E(u_{0n})\zeta(x,0)\,\mathrm{d}x .
\end{align*}
取 $E=E_l^k$，其中 $E_l^k\in C^2([0,\infty))$ 为凸函数、$E_l^k\to E_l$ 在 $[0,\infty)$ 上局部一致且 $|(E_l^k)'|\leqslant1$，相应的 $F_l^k$ 由 $(F_l^k)'=(E_l^k)'\phi'$ 与 $F_l^k(l)=0$ 确定，则 $F_l^k\to F_l$ 在 $[0,\infty)$ 上局部一致；对 $E_l^k$ 用上式并令 $k\to\infty$，即得\eqref{align:ch07-7.77}。故 $y_n$ 是 (Pn) 的熵解；因 $\phi\in C^\infty([0,\infty))$ 局部 Lipschitz 连续，由 Kružkov 唯一性定理（见文献~\cite{ref91}）得 $y_n=u_n$。结论得证。
\end{proof}

% SUPPLEMENT: parent=lemma:ch07-7.6.1; kind=motivation; 说明 (V$n^\varepsilon$) 作为人工粘性正则化的动机与引理在\cref{theorem:ch07-7.3.3}证明中的枢纽地位，逐项说明 (A2')、$\phi''<0$、$\phi'>0$ 与 $u_{0n}$ 的\eqref{align:ch07-7.44}--\eqref{align:ch07-7.47}的作用，并给出 $\phi'$ 取零值、$u_{0n}\notin L^\infty$、$\phi$ 非局部 Lipschitz 三种边界情形
\begin{remark}\phantomsection\label{remark:ch07-Vn逼近}
(i) 动机。(V$n^\varepsilon$) 把 (Pn) 的对流项 $\partial_x[\phi(y)]$ 换成人工粘性 $\varepsilon\partial_{xx}[\phi(y)]$，其解光滑，并在 $\varepsilon\to0^+$ 时收敛于 (Pn) 的熵解。\cref{lemma:ch07-7.6.1}的作用是把\cref{section:ch07-逼近问题}中为 (P$n^\varepsilon$) 建立的紧性结论移植到 (V$n^\varepsilon$)：一旦 $u_{n_k}$ 可由光滑解 $y_{n_k}^{\varepsilon_j}$ 逼近，就可对 $y_{n_k}^{\varepsilon_j}$ 使用\cref{lemma:ch07-7.6.2}的时间单调估计，再令 $\varepsilon_j\to0^+$；这是证明\cref{theorem:ch07-7.3.3}的通道。

(ii) 各假设的作用。(A2') 中 $\phi\in C^\infty([0,\infty))$ 与 $\lim_{y\to\infty}\phi(y)/y=0$ 使 (V$n^\varepsilon$) 的系数光滑且在 $y_n^\varepsilon$ 的值域上一致有界，从而存在性、唯一性与 $L^\infty$ 估计\eqref{align:ch07-7.49}成立。$\phi'>0$ 在 $(0,\infty)$ 是扩散项 $\varepsilon\partial_{xx}[\phi(y)]$ 的非退化条件：它使 $\phi$ 严格增，从而 (V$n^\varepsilon$) 为拟线性一致抛物方程，且 $\phi(y_n^\varepsilon)$ 的 $L^1$ 界可控制 $y_n^\varepsilon$ 的 $L^1$ 界。$\phi''<0$ 与 $\alpha\geqslant-1$ 使 $g=(\alpha\phi+\beta)/\phi'$ 满足 $g\geqslant0$ 与 $g'\geqslant\alpha+1$（见\cref{lemma:ch07-7.6.2}）。$u_{0n}$ 的\eqref{align:ch07-7.44}--\eqref{align:ch07-7.45}给出与 $n$ 无关的 $L^1$ 与 $L^\infty$ 界，这是\eqref{align:ch07-7.66}与\eqref{align:ch07-7.69}中估计关于 $n$ 一致的来源；$u_{0n}^\varepsilon$ 的\eqref{align:ch07-7.46}--\eqref{align:ch07-7.47}使初值在 $L^1$ 中强收敛、在 $L^\infty$ 中弱*收敛，故熵不等式中的初值项可取极限。$u_n$ 的唯一性使极限 $y_n$ 可与 $u_n$ 直接比较（Kružkov 定理），否则只能断言 $y_n$ 是某个熵解。

(iii) 边界与反例。若 $\phi'$ 在某个区间上为零（$\phi$ 有平坦段），则 $\phi$ 不严格增，$g$ 无定义，$\phi(y_n^\varepsilon)$ 不能控制 $y_n^\varepsilon$，\eqref{align:ch07-7.165}的 $L^1_{\mathrm{loc}}$ 收敛不能由 $\phi(y_n^{\varepsilon_j})$ 的收敛推出；故 $\phi'>0$ 不可省。若去掉 $u_{0n}\in L^\infty(\mathbb{R})$，则\eqref{align:ch07-7.49}型的 $L^\infty$ 界失效，弱*收敛子列仍存在，但 $\{y_n^\varepsilon\}$ 的 $L^1_{\mathrm{loc}}$ 紧性需另行论证。若 $\phi$ 仅连续而不局部 Lipschitz，(Pn) 的熵解不唯一，结论须改为“$y_n$ 是 (Pn) 的某一熵解”。
\end{remark}

\begin{lemma}\label{lemma:ch07-7.6.2}
设 (A0) 与 (A2') 成立，$\phi(0)=0$，$\phi''<0$ 与 $\phi'>0$ 在 $(0,\infty)$。设 $n\in\mathbb{N}$、$\varepsilon>0$，$u_{0n}^\varepsilon\in C_c^\infty(\mathbb{R})$ 非负，$y_n^\varepsilon\in L^\infty(S)\cap L^\infty(0,T;L^1(\mathbb{R}))$ 为 (V$n^\varepsilon$) 在 $S$ 中的解。则对一切 $t\in(0,T)$，函数 $x\mapsto\phi(y_n^\varepsilon(x,t))$ 属于 $W^{1,1}_{\mathrm{loc}}(\mathbb{R})$，且
\begin{subequations}
\begin{align}
\pm\frac{\partial}{\partial t}\left[\frac{\alpha\phi(y_n^\varepsilon(\cdot,t))+\beta}{t^\alpha}\right]\leqslant0\quad\text{在 }\mathbb{R}\text{ 中几乎处处（若 }\pm\alpha>0\text{）,} \label{align:ch07-7.166a}
\end{align}
\begin{align}
\frac{\partial}{\partial t}[\phi(y_n^\varepsilon(\cdot,t))-\beta\log t]\leqslant0\quad\text{在 }\mathbb{R}\text{ 中几乎处处（若 }\alpha=0\text{）.} \label{align:ch07-7.166b}
\end{align}
\end{subequations}
\end{lemma}

% SUPPLEMENT: parent=lemma:ch07-7.6.2; kind=proof; 补出 $\partial_tz_n^\varepsilon$ 的完整计算（含 $A[\alpha\phi(y)+\beta]=\alpha\partial_ty_n^\varepsilon$ 一步）、由 (A2') 与 $\alpha\geqslant-1$ 得 $g'\geqslant\alpha+1$ 的推导、$z_n^\varepsilon(\cdot,0)\leqslant0$ 与零阶系数非正时比较原理的适用，以及由 $z_n^\varepsilon\leqslant0$ 求导得\eqref{align:ch07-7.166a}--\eqref{align:ch07-7.166b}的计算
\begin{proof}
为简化记号置 $A\triangleq\varepsilon\partial_{xx}-\partial_x$，于是 (V$n^\varepsilon$) 的第一式即
\begin{align*}
\partial_ty_n^\varepsilon=A[\phi(y_n^\varepsilon)]\quad\text{在 }S .
\end{align*}
置
\begin{align*}
z_n^\varepsilon\triangleq t\partial_ty_n^\varepsilon-g(y_n^\varepsilon),\qquad\text{其中 }
g(y)\triangleq\frac{\alpha\phi(y)+\beta}{\phi'(y)}\quad(y\in[0,\infty)).
\end{align*}
因 $\phi'>0$，$g$ 在 $[0,\infty)$ 上有定义；因 $\phi''<0$ 且 $\phi''[\alpha\phi+\beta]\leqslant-(\phi')^2<0$，有 $\alpha\phi+\beta>0$ 在 $[0,\infty)$，故 $g\geqslant0$，且 $g'\geqslant\alpha+1\geqslant0$：事实上
\begin{align*}
g'=\frac{\alpha(\phi')^2-(\alpha\phi+\beta)\phi''}{(\phi')^2}=\alpha-\frac{(\alpha\phi+\beta)\phi''}{(\phi')^2}\geqslant\alpha+1,
\end{align*}
最后一步用 $(\alpha\phi+\beta)\phi''/(\phi')^2\leqslant-1$（即 (A2')）与 $\alpha\geqslant-1$。

由 (A2') 与直接计算，
\begin{align*}
\partial_tz_n^\varepsilon&=\partial_ty_n^\varepsilon+t\partial_{tt}y_n^\varepsilon-g'(y_n^\varepsilon)\partial_ty_n^\varepsilon
=A[t\partial_t\phi(y_n^\varepsilon)]+[1-g'(y_n^\varepsilon)]\partial_ty_n^\varepsilon\\
&=A[\phi'(y_n^\varepsilon)(z_n^\varepsilon+g(y_n^\varepsilon))]+[1-g'(y_n^\varepsilon)]\frac{z_n^\varepsilon+g(y_n^\varepsilon)}{t}\\
&=A[\phi'(y_n^\varepsilon)z_n^\varepsilon]+A[\alpha\phi(y_n^\varepsilon)+\beta]+\frac{[\alpha+1-g'(y_n^\varepsilon)][z_n^\varepsilon+g(y_n^\varepsilon)]}{t}\\
&=A[\phi'(y_n^\varepsilon)z_n^\varepsilon]+\alpha\partial_ty_n^\varepsilon+\frac{[\alpha+1-g'(y_n^\varepsilon)][z_n^\varepsilon+g(y_n^\varepsilon)]}{t}\\
&=A[\phi'(y_n^\varepsilon)z_n^\varepsilon]+\frac{[\alpha+1-g'(y_n^\varepsilon)][z_n^\varepsilon+g(y_n^\varepsilon)]}{t}\quad\text{在 }S,
\end{align*}
其中第三、五步分别用 $t\partial_ty_n^\varepsilon=z_n^\varepsilon+g(y_n^\varepsilon)$ 与 $t\phi'(y_n^\varepsilon)\partial_ty_n^\varepsilon=\phi'(y_n^\varepsilon)z_n^\varepsilon+\alpha\phi(y_n^\varepsilon)+\beta$（即 $\phi'g=\alpha\phi+\beta$），第四步用 $A[\alpha\phi(y)+\beta]=\alpha A[\phi(y)]=\alpha\partial_ty_n^\varepsilon$。因 $g\geqslant0$ 且 $\alpha+1-g'\leqslant0$，上式给出
\begin{align*}
\partial_tz_n^\varepsilon\leqslant A[\phi'(y_n^\varepsilon)z_n^\varepsilon]+\frac{\alpha+1-g'(y_n^\varepsilon)}{t}z_n^\varepsilon\quad\text{在 }S .
\end{align*}
这是关于 $z_n^\varepsilon$ 的线性方程：二阶项系数为 $\varepsilon\phi'(y_n^\varepsilon)>0$，零阶系数为 $[\alpha+1-g'(y_n^\varepsilon)]/t\leqslant0$。因 $z_n^\varepsilon(\cdot,0)=-g(u_{0n}^\varepsilon)\leqslant0$ 在 $\mathbb{R}$（用到 $g\geqslant0$ 与 $u_{0n}^\varepsilon\geqslant0$），由比较原理（常数 $0$ 是该线性方程的上解）得 $z_n^\varepsilon\leqslant0$ 在 $S$ 中，即
\begin{align*}
t\partial_ty_n^\varepsilon(\cdot,t)\leqslant g(y_n^\varepsilon)(\cdot,t)\quad\text{在 }\mathbb{R}\text{ 中几乎处处，对一切 }t\in(0,T).
\end{align*}
由 $\phi'>0$ 得 $t\phi'(y_n^\varepsilon)\partial_ty_n^\varepsilon\leqslant\alpha\phi(y_n^\varepsilon)+\beta$，而
\begin{align*}
\frac{\partial}{\partial t}\left[\frac{\alpha\phi(y_n^\varepsilon)+\beta}{t^\alpha}\right]=\frac{\alpha\left[t\phi'(y_n^\varepsilon)\partial_ty_n^\varepsilon-(\alpha\phi(y_n^\varepsilon)+\beta)\right]}{t^{\alpha+1}},
\end{align*}
故 $\alpha>0$ 时上式左端 $\leqslant0$，$\alpha<0$ 时 $\geqslant0$，即\eqref{align:ch07-7.166a}。若 $\alpha=0$，则
\begin{align*}
\frac{\partial}{\partial t}[\phi(y_n^\varepsilon)-\beta\log t]=\phi'(y_n^\varepsilon)\partial_ty_n^\varepsilon-\frac{\beta}{t}\leqslant\frac{\alpha\phi(y_n^\varepsilon)+\beta}{t}-\frac{\beta}{t}=0,
\end{align*}
即\eqref{align:ch07-7.166b}。结论得证。
\end{proof}

% SUPPLEMENT: parent=lemma:ch07-7.6.2; kind=motivation; 说明该时间单调估计作为 Bénilan--Crandall 型先验估计的动机与其在\cref{theorem:ch07-7.3.3}(ii) 中的作用，逐项说明 $\alpha\phi+\beta>0$、$g\geqslant0$、$g'\geqslant\alpha+1$ 的必要性，并给出 $\alpha=0$、$\alpha=-1$、$\alpha<-1$ 三种边界情形
\begin{remark}\phantomsection\label{remark:ch07-时间单调性}
(i) 动机。\eqref{align:ch07-7.166a}--\eqref{align:ch07-7.166b}是 Bénilan--Crandall 型（Aronson--Bénilan）时间正则化估计：它把 (V$n^\varepsilon$) 的解在正时刻的增长率限制为 $t^{-\alpha}$ 阶，且界与空间变量 $x$ 无关。该估计是\cref{theorem:ch07-7.3.3}(ii) 中“$\phi$ 无界时 $u=u_r\in L^\infty(\mathbb{R}\times(\tau,T))$”的来源：沿时间积分\eqref{align:ch07-7.166a}得\eqref{align:ch07-7.173}，再由\eqref{align:ch07-7.173}与 $Y_n$ 的 $L^1$ 界得\eqref{align:ch07-7.174}。

(ii) 各假设的作用。(A0) 与 $u_{0n}^\varepsilon\geqslant0$ 使 $z_n^\varepsilon(\cdot,0)=-g(u_{0n}^\varepsilon)\leqslant0$，这是比较原理的初值条件；$g\geqslant0$ 又等价于 $\alpha\phi+\beta>0$，即情形 (a)，故该估计只在情形 (a) 可用。$\phi'>0$ 使 $g$ 有定义，并使 $z_n^\varepsilon$ 的方程中二阶项系数 $\varepsilon\phi'$ 为正，比较原理才可用。$\phi''<0$ 与 $\alpha\geqslant-1$ 合起来给出 $g'\geqslant\alpha+1$，从而零阶系数 $[\alpha+1-g']/t$ 非正：若该系数可取正值，$z_n^\varepsilon$ 的正部会被放大，$z_n^\varepsilon\leqslant0$ 不再成立。

(iii) 边界。$\alpha=0$ 时\eqref{align:ch07-7.166a}两端都是常数 $\beta$，不等式不含信息，故须单列\eqref{align:ch07-7.166b}。$\alpha=-1$ 时 $t^{\alpha+1}-t_1^{\alpha+1}=0$，\eqref{align:ch07-7.173}的第一项消失，\eqref{align:ch07-7.174}须由 $t\mapsto t(\beta-\phi(y_n^\varepsilon))$ 的单调性直接得到（见\cref{theorem:ch07-7.3.3}(ii) 的证明）。若 $\alpha<-1$，则\eqref{align:ch07-7.174}的常数 $(\alpha+1)\|u_0\|_{\mathscr{R}_f(\mathbb{R})}/t+\beta$ 当 $t$ 充分小时为负，与 $\phi\geqslant0$ 矛盾，故 $\alpha\geqslant-1$ 不可放宽。
\end{remark}

\begin{lemma}\label{lemma:ch07-7.6.3}
设 (A0)--(A1) 成立，$u$ 为\cref{definition:ch07-构造解}意义下的所构造解，$\{u_{n_k}\}$ 为\cref{theorem:ch07-7.5.1}与\cref{remark:ch07-7.5.1}证明中所取的子列（其弱*极限给出 $u=u_r\,\mathrm{d}x+u_s$）。则对几乎所有 $t\in(0,T)$ 与每个 $x_0\in\operatorname{supp}u_s(\cdot,t)$，存在序列 $\{x_p\}\subseteq\mathbb{R}$ 与子列 $\{u_{n_p}\}\subseteq\{u_{n_k}\}$，使 $x_p\to x_0$ 且 $u_{n_p}(x_p,t)\to\infty$（$p\to\infty$）。
\end{lemma}

% SUPPLEMENT: parent=lemma:ch07-7.6.3; kind=proof; 补出\eqref{align:ch07-7.167}中零测集与子列的取法、$u_s(\cdot,t)$ 的绝对连续性导致矛盾的完整论证（含 $f_t\geqslant0$ 的来源与弱*极限在 $C_c(I_\delta(x_0))$ 上的识别）、以及对角线选取使同一个 $\{u_{n_p}\}$ 同时满足 $x_p\in I_{1/p}(x_0)$ 与 $u_{n_p}(x_p,t)\geqslant p$
\begin{proof}
由\cref{remark:ch07-7.5.1}，存在零测集 $N\subseteq(0,T)$ 与 $\{u_{n_k}\}$ 的子列（不重新标号），使
\begin{align}
u_{n_k}(\cdot,t)\stackrel{*}{\rightharpoonup}u_r(\cdot,t)+u_s(\cdot,t)\quad\text{在 }\mathscr{R}_f(I_L) \label{align:ch07-7.167}
\end{align}
对所有 $t\in(0,T)\setminus N$ 与 $L>0$ 成立（见\eqref{align:ch07-7.107}取 $U(y)=y$ 及\eqref{align:ch07-7.109a}--\eqref{align:ch07-7.109b}）。

固定 $t\in(0,T)\setminus N$ 并设 $x_0\in\operatorname{supp}u_s(\cdot,t)$。断言不存在邻域 $I_\delta(x_0)$ 使 $\{u_{n_k}(\cdot,t)\}$ 含于 $L^\infty(I_\delta(x_0))$ 的有界子集。事实上，若存在这样的邻域，则由 Banach--Alaoglu 定理，存在子列（不重新标号）与 $f_t\in L^\infty(I_\delta(x_0))$，使 $u_{n_k}(\cdot,t)\stackrel{*}{\rightharpoonup}f_t$ 在 $L^\infty(I_\delta(x_0))$；因 $u_{n_k}(\cdot,t)\geqslant0$ 几乎处处，取非负试验函数即得 $f_t\geqslant0$。又 $C_c(I_\delta(x_0))\subseteq L^1(I_\delta(x_0))\subseteq(L^\infty(I_\delta(x_0)))^*$，故同一子列在 $\mathscr{R}_f(I_\delta(x_0))$ 中弱*收敛于 $f_t\,\mathrm{d}x$；与\eqref{align:ch07-7.167}比较（弱*极限唯一）得 $u_r(\cdot,t)+u_s(\cdot,t)=f_t\,\mathrm{d}x$ 在 $I_\delta(x_0)$ 上，即 $u_s(\cdot,t)$ 在 $I_\delta(x_0)$ 上为零测度。这与 $x_0\in\operatorname{supp}u_s(\cdot,t)$（即每个邻域上 $|u_s(\cdot,t)|$ 的质量为正）矛盾。

于是 $\sup_{n_k\in\mathbb{N}}\|u_{n_k}(\cdot,t)\|_{L^\infty(I_{1/p}(x_0))}=\infty$ 对一切 $p\in\mathbb{N}$ 成立。由此用对角线方法选取子列 $\{u_{n_{k_p}}\}\subseteq\{u_{n_k}\}$ 与点 $x_p\in I_{1/p}(x_0)$：在第 $p$ 步，因 $\|u_{n}(\cdot,t)\|_{L^\infty(I_{1/p}(x_0))}$ 关于 $n$ 无界，存在 $n=n_{k_p}$（指标大于前一步所取）与 $L^1(I_{1/p}(x_0))$ 中函数 $u_{n_{k_p}}(\cdot,t)$ 的 Lebesgue 点 $x_p\in I_{1/p}(x_0)$，使 $u_{n_{k_p}}(x_p,t)\geqslant p$。则 $x_p\to x_0$ 且 $u_{n_p}(x_p,t)\geqslant p\to\infty$。结论得证。
\end{proof}

% SUPPLEMENT: parent=lemma:ch07-7.6.3; kind=motivation; 说明该“支撑点处密度爆破”引理的动机与其在\cref{theorem:ch07-7.3.3}(i) 中导出\eqref{align:ch07-7.171}的作用，逐项说明 (A0)--(A1)、所构造解、$x_0\in\operatorname{supp}u_s(\cdot,t)$ 的必要性，并给出 $x_0\notin\operatorname{supp}u_s(\cdot,t)$ 与正则时刻两种边界情形
\begin{remark}\phantomsection\label{remark:ch07-支撑爆破}
(i) 动机。奇异部分 $u_s(\cdot,t)$ 的支撑点处，解的密度必沿某一侧无界。\cref{lemma:ch07-7.6.3}给出的正是这一现象在逼近序列上的形式：存在趋于 $x_0$ 的点 $x_p$ 与相应的子列，使 $u_{n_p}(x_p,t)$ 任意大。它是\cref{theorem:ch07-7.3.3}(i) 中\eqref{align:ch07-7.171}的出发点：把\eqref{align:ch07-7.170}用于以 $x_p$ 为左端点的试验函数，即可把 $Z(u_{n_p}(x_p,t))$ 换成 $Z(u_r(\cdot,t))$ 在 $x_0$ 右侧的上极限，从而说明奇异质量的支撑点必为密度的爆破点。

(ii) 各假设的作用。(A0)--(A1) 使\cref{theorem:ch07-7.5.1}与\cref{remark:ch07-7.5.1}可用，从而存在公共零测集 $N$ 使\eqref{align:ch07-7.107}（取 $U(y)=y$）对一切 $t\in(0,T)\setminus N$ 成立，这是\eqref{align:ch07-7.167}的来源。$u$ 为所构造解保证 $\{u_{n_k}\}$ 的弱*极限恰为 $u_r\,\mathrm{d}x+u_s$，故 $u_s(\cdot,t)$ 的支撑与 $\{u_{n_k}(\cdot,t)\}$ 的局部无界性可以比较；若把 $u$ 换成任意解，\eqref{align:ch07-7.167}不再成立。$x_0\in\operatorname{supp}u_s(\cdot,t)$ 是非退化条件：$u_s(\cdot,t)$ 的每个邻域上质量为正，故不存在使 $\{u_{n_k}(\cdot,t)\}$ 局部 $L^\infty$ 有界的邻域。

(iii) 边界。结论只对 $t\in(0,T)\setminus N$ 成立；$x_p$ 与子列同时依赖 $p$，不能对固定的 $n$ 取得无界值。若 $x_0\notin\operatorname{supp}u_s(\cdot,t)$，则 $u_s(\cdot,t)$ 在某个 $I_\delta(x_0)$ 上为零，由 $u_r\in L^1_{\mathrm{loc}}$ 与\eqref{align:ch07-7.167}知 $\{u_{n_k}(\cdot,t)\}$ 在 $I_{\delta/2}(x_0)$ 上 $L^\infty$ 有界，断言不成立。若 $u_s(\cdot,t)=0$（正则时刻），则 $\operatorname{supp}u_s(\cdot,t)=\varnothing$，断言空洞。
\end{remark}

% SUPPLEMENT: parent=theorem:ch07-7.3.3; kind=proof; 补出 (i) 中的三处实质步骤：由\eqref{align:ch07-7.169}乘 $\rho/g(y)$ 得\eqref{align:ch07-7.170}的乘子计算（含 $Z'=\phi'/g$ 与 $g'\geqslant0$ 的符号）、由\eqref{align:ch07-7.170}经 $\rho_{p,m}$ 与 Lebesgue 点得\eqref{align:ch07-7.171}--\eqref{align:ch07-7.172}、以及相容条件 (C$^\pm$) 的逐项验证
\begin{proof}[\cref{theorem:ch07-7.3.3}的证明]
如前所述，只需就 $\phi''<0$ 与 $\phi'>0$ 在 $(0,\infty)$ 的情形证明。

(i) 设 $\phi$ 有界，$\{u_{n_k}\}$ 如\cref{lemma:ch07-7.6.3}证明中。由\cref{lemma:ch07-7.6.1}，对每个 $n_k\in\mathbb{N}$，存在 (V$n_k^\varepsilon$) 解的子列 $\{y_{n_k}^{\varepsilon_j}\}$，使对几乎所有 $t\in(0,T)$，
\begin{align}
y_{n_k}^{\varepsilon_j}(\cdot,t)\to u_{n_k}(\cdot,t)\quad\text{在 }L^1_{\mathrm{loc}}(\mathbb{R})\ \text{当 }j\to\infty. \label{align:ch07-7.168}
\end{align}
（对可数个 $n_k$ 用对角线论证，可取公共的 $\varepsilon_j\to0^+$。）另一方面，由\cref{lemma:ch07-7.6.2}的证明，对任意 $t\in(0,T)$，有
\begin{align}
\varepsilon_j\partial_{xx}\phi(y_{n_k}^{\varepsilon_j}(\cdot,t))-\partial_x\phi(y_{n_k}^{\varepsilon_j}(\cdot,t))
=\partial_t y_{n_k}^{\varepsilon_j}(\cdot,t)\leqslant\frac{g(y_{n_k}^{\varepsilon_j}(\cdot,t))}{t}\quad\text{在 }\mathbb{R}, \label{align:ch07-7.169}
\end{align}
其中 $g(y)\triangleq(\alpha\phi(y)+\beta)/\phi'(y)>0$。

固定 $a,b\in\mathbb{R}$，$a<b$，设 $\rho\in C_c^1(a,b)$，$\rho\geqslant0$，并置
\begin{align*}
Z(y)\triangleq\int_0^y\frac{[\phi'(z)]^2}{\alpha\phi(z)+\beta}\,\mathrm{d}z\quad(y\in[0,\infty)).
\end{align*}
则 $Z$ 递增，且 $0\leqslant Z(y)\leqslant\phi'(0)$ 对一切 $y\geqslant0$：事实上由 (A2') 得 $[\phi'(z)]^2/(\alpha\phi(z)+\beta)\leqslant-\phi''(z)$，故 $Z(y)\leqslant\phi'(0)-\phi'(y)\leqslant\phi'(0)$。将\eqref{align:ch07-7.169}乘以 $\rho/g(y_{n_k}^{\varepsilon_j}(\cdot,t))$ 并在 $(a,b)$ 上积分。由 $Z'=\phi'/g$（即 $Z'(y)=\phi'(y)/g(y)$）得 $\partial_xZ(y)=\phi'(y)\partial_xy/g(y)$，分部积分给出
\begin{align*}
&\int_a^bZ(y_{n_k}^{\varepsilon_j})(x,t)[\varepsilon_j\rho''(x)+\rho'(x)]\,\mathrm{d}x\\
&\qquad=\frac1t\int_a^b\rho(x)\,\mathrm{d}x
-\varepsilon_j\int_a^b\frac{\phi'(y_{n_k}^{\varepsilon_j})g'(y_{n_k}^{\varepsilon_j})}{[g(y_{n_k}^{\varepsilon_j})]^2}[\partial_x y_{n_k}^{\varepsilon_j}]^2(x,t)\rho(x)\,\mathrm{d}x\\
&\qquad\leqslant\frac1t\int_a^b\rho(x)\,\mathrm{d}x,
\end{align*}
最后一步用 $\rho\geqslant0$、$\phi'>0$ 与 $g'\geqslant0$；而
\begin{align*}
g'=\alpha-\frac{(\alpha\phi+\beta)\phi''}{(\phi')^2}\geqslant\alpha+1\geqslant0
\end{align*}
由 (A2') 与 $\alpha\geqslant-1$ 得到。由\eqref{align:ch07-7.168}，必要时再取子列，$y_{n_k}^{\varepsilon_j}(x,t)\to u_{n_k}(x,t)$ 对几乎所有 $x\in(a,b)$ 成立；因 $0\leqslant Z\leqslant\phi'(0)$ 而 $Z$ 连续，由控制收敛定理可令 $j\to\infty$，得
\begin{align}
\int_a^bZ(u_{n_k}(x,t))\rho'(x)\,\mathrm{d}x\leqslant\frac1t\int_a^b\rho(x)\,\mathrm{d}x. \label{align:ch07-7.170}
\end{align}

对几乎所有 $t\in(0,T)$，设 $x_0\in\operatorname{supp}u_s(\cdot,t)$。证明
\begin{align}
\operatorname*{ess\,lim}_{x\to x_0^+}u_r(x,t)=\infty. \label{align:ch07-7.171}
\end{align}
设 $\{x_p\}$ 与 $\{u_{n_p}\}\subseteq\{u_{n_k}\}$ 由\cref{lemma:ch07-7.6.3}给出，并固定 $b>x_0$。因 $x_p\to x_0$（$p\to\infty$），存在 $\bar p\in\mathbb{N}$ 使 $b>x_p$ 对一切 $p>\bar p$ 成立。对这样的 $p$ 与充分大的 $m\in\mathbb{N}$ 置
\begin{align*}
\rho_{p,m}(x)\triangleq m(x-x_p)\chi_{[x_p,x_p+1/m]}(x)+\chi_{(x_p+1/m,b-1/m)}(x)+m(b-x)\chi_{[b-1/m,b]}(x).
\end{align*}
在\eqref{align:ch07-7.170}中取 $a=x_p$ 与 $\rho=\rho_{p,m}$：因 $\rho_{p,m}'=m$ 在 $(x_p,x_p+1/m)$ 上、$\rho_{p,m}'=0$ 在 $(x_p+1/m,b-1/m)$ 上、$\rho_{p,m}'=-m$ 在 $(b-1/m,b)$ 上，得
\begin{align*}
m\int_{x_p}^{x_p+1/m}Z(u_{n_p}(x,t))\,\mathrm{d}x-m\int_{b-1/m}^bZ(u_{n_p}(x,t))\,\mathrm{d}x\leqslant\frac1t\int_{x_p}^b\rho_{p,m}(x)\,\mathrm{d}x,
\end{align*}
从而令 $m\to\infty$（左端两个平均值分别收敛于 $Z(u_{n_p}(x_p,t))$ 与 $Z(u_{n_p}(b,t))$，右端收敛于 $(b-x_p)/t$；这里取 $x_p$ 为 $Z(u_{n_p}(\cdot,t))$ 的 Lebesgue 点、$b$ 为 $Z(u_{n_p}(\cdot,t))$ 的 Lebesgue 点，由\cref{lemma:ch07-7.6.3}的证明中 $x_p$ 可在 $\{u_{n_p}(\cdot,t)\geqslant p\}\cap I_{1/p}(x_0)$ 的密度点中选取，这不妨碍 $u_{n_p}(x_p,t)\geqslant p$），对所有 $p>\bar p$ 与几乎所有 $b>x_0$ 有
\begin{align*}
Z(u_{n_p}(x_p,t))\leqslant Z(u_{n_p}(b,t))+\frac{b-x_p}{t}.
\end{align*}
由\cref{lemma:ch07-7.6.3}的证明，$u_{n_p}(x_p,t)\geqslant p$。令 $p\to\infty$ 并利用\eqref{align:ch07-7.143}（$u_{n_k}\to u_r$ 在 $S$ 中几乎处处，故 $Z(u_{n_p}(b,t))\to Z(u_r(b,t))$ 对几乎所有 $b$），得对几乎所有 $b>x_0$
\begin{align}
\lim_{y\to\infty}Z(y)\leqslant Z(u_r(b,t))+\frac{b-x_0}{t}, \label{align:ch07-7.172}
\end{align}
从而
\begin{align*}
\lim_{y\to\infty}Z(y)\leqslant\operatorname*{ess\,lim\,inf}_{b\to x_0^+}Z(u_r(b,t)).
\end{align*}
因 $Z$ 严格递增且 $\lim_{y\to\infty}Z(y)=\sup_{y\geqslant0}Z(y)=\phi'(0)$，上式蕴含 $Z(u_r(b,t))\to\sup Z$（$b\to x_0^+$），故 $u_r(b,t)\to\infty$，即\eqref{align:ch07-7.171}。

现证 $u$ 在 $[0,\tau]$ 上于 $x_0\in\operatorname{supp}u_s(\cdot,\tau)$ 处满足相容条件 (C$^\pm$)。关于 (C$^-$)：由 $\phi$ 的递增性，若 $u_r>l$ 则 $H_-(u_r-l)=0$；若 $u_r\leqslant l$ 则
\begin{align*}
H_-(u_r(x,t)-l)[\phi(u_r(x,t))-\phi(l)]=\phi(l)-\phi(u_r(x,t))\geqslant0,
\end{align*}
故被积函数非负，其单侧极限（由\cref{lemma:ch07-7.2.3}存在）非负，即 (C$^-$) 成立。关于 (C$^+$)：由\eqref{align:ch07-7.171}，对一切 $l\in[0,\infty)$ 与几乎所有 $t\in(0,\tau)$ 有 $\operatorname*{ess\,lim}_{x\to x_0^+}H_-(u_r(x,t)-l)=0$（由 $u_s(\cdot,t)$ 关于 $t$ 的单调性\eqref{align:ch07-7.15}，$x_0\in\operatorname{supp}u_s(\cdot,t)$ 对这样的 $t$ 成立）；于是由控制收敛定理，
\begin{align*}
\operatorname*{ess\,lim}_{x\to x_0^+}\int_0^\tau H_-(u_r(x,t)-l)[\phi(u_r(x,t))-\phi(l)]\beta(t)\,\mathrm{d}t=0
\end{align*}
对一切 $l\in[0,\infty)$ 与 $\beta\in C_c^1(0,\tau)$、$\beta\geqslant0$ 成立，即 (C$^+$) 成立。

% SUPPLEMENT: parent=theorem:ch07-7.3.3; kind=proof; 补出 (ii) 中由\eqref{align:ch07-7.166a}沿时间积分得\eqref{align:ch07-7.173}的逐项计算（含 $\pm\alpha$ 与 $1/\alpha$ 的符号相消）、$\alpha=-1$ 的极限情形、$Y_n\in W^{1,1}$ 与 $0\leqslant Y_n\leqslant\|u_0\|_{\mathscr{R}_f(\mathbb{R})}$ 的推导、由\eqref{align:ch07-7.173}取 $t_1\to0^+$ 得\eqref{align:ch07-7.174}（含 $\alpha\geqslant-1$ 的作用）、无界 $\phi$ 时 $\alpha\geqslant0$ 的理由，以及最后取 $n_k\to\infty$ 的依据
(ii) 设 $\phi$ 无界。由\eqref{align:ch07-7.166a}，对一切 $0<t_1\leqslant t\leqslant T$ 与 $x\in\mathbb{R}$，
\begin{align*}
\int_{t_1}^t\phi(y_{n_k}^{\varepsilon_j})(x,s)\,\mathrm{d}s
&=\int_{t_1}^t\frac{\alpha\phi(y_{n_k}^{\varepsilon_j})(x,s)+\beta}{s^\alpha}s^\alpha\,\mathrm{d}s-\frac{\beta}{\alpha}(t-t_1)\\
&\geqslant\frac{\alpha\phi(y_{n_k}^{\varepsilon_j})(x,t)+\beta}{\alpha t^\alpha}\cdot\frac{t^{\alpha+1}-t_1^{\alpha+1}}{\alpha+1}-\frac{\beta}{\alpha}(t-t_1)\quad\text{若 }\alpha\neq0,\\
\int_{t_1}^t\phi(y_{n_k}^{\varepsilon_j})(x,s)\,\mathrm{d}s
&=\int_{t_1}^t[\phi(y_{n_k}^{\varepsilon_j})(x,s)-\beta\log s]\,\mathrm{d}s+\beta\int_{t_1}^t\log s\,\mathrm{d}s\\
&\geqslant[\phi(y_{n_k}^{\varepsilon_j})(x,t)-\beta\log t](t-t_1)+\beta[t\log t-t]-\beta[t_1\log t_1-t_1]\quad\text{若 }\alpha=0,
\end{align*}
其中 $\{y_{n_k}^{\varepsilon_j}\}$ 为\cref{lemma:ch07-7.6.1}证明中所用子列。令 $\varepsilon_j\to0^+$，由\eqref{align:ch07-7.165}得，对几乎所有 $t\in(t_1,T)$ 与几乎所有 $x\in\mathbb{R}$，
\begin{align}
\int_{t_1}^t\phi(u_{n_k})(x,s)\,\mathrm{d}s\geqslant
\begin{cases}
\dfrac{\alpha\phi(u_{n_k})(x,t)+\beta}{\alpha t^\alpha}\dfrac{t^{\alpha+1}-t_1^{\alpha+1}}{\alpha+1}-\dfrac{\beta}{\alpha}(t-t_1)&\text{若 }\alpha\neq0,\\[2mm]
[\phi(u_{n_k})(x,t)-\beta](t-t_1)+\beta t_1\log\dfrac{t}{t_1}&\text{若 }\alpha=0,
\end{cases} 
\label{align:ch07-7.173}
\end{align}
（$\alpha=0$ 一行的整理用 $-\beta\log t(t-t_1)+\beta(t\log t-t)-\beta(t_1\log t_1-t_1)=-\beta(t-t_1)+\beta t_1\log(t/t_1)$）。

另一方面，如\cref{proposition:ch07-7.2.1}的证明那样，由弱形式\eqref{align:ch07-7.78}（对 (Pn) 的熵解 $u_{n_k}$）得：对每个 $t\in(0,T]$，映射 $x\mapsto Y_n(x)\triangleq\int_0^t\phi(u_{n_k}(x,s))\,\mathrm{d}s$ 属于 $W^{1,1}(\mathbb{R})$，其弱导数 $Y_n'(\cdot)=u_{0n_k}-u_{n_k}(\cdot,t)$ 在 $\mathbb{R}$ 中几乎处处成立。故
\begin{align*}
0\leqslant Y_n\leqslant\|u_{0n_k}\|_{L^1(\mathbb{R})}\leqslant\|u_0\|_{\mathscr{R}_f(\mathbb{R})},
\end{align*}
其中第一、二个不等号分别用 $\phi\geqslant0$（$u_{n_k}\geqslant0$ 与 $\phi(0)=0$、$\phi'>0$）与 $Y_n(\pm\infty)=0$、$\operatorname{Var}(Y_n)\leqslant\|u_{0n_k}-u_{n_k}(\cdot,t)\|_{L^1(\mathbb{R})}\leqslant2\|u_0\|_{\mathscr{R}_f(\mathbb{R})}$（即 $Y_n\geqslant0$ 且两端为 $0$ 时 $\sup Y_n\leqslant\operatorname{Var}(Y_n)/2$）。

结合此不等式与\eqref{align:ch07-7.173}并令 $t_1\to0^+$，得
\begin{align}
\phi(u_{n_k})(x,t)\leqslant\frac{(\alpha+1)\|u_0\|_{\mathscr{R}_f(\mathbb{R})}}{t}+\beta
\quad\text{对几乎所有 }t\in(0,T)\text{ 与几乎所有 }x\in\mathbb{R} \label{align:ch07-7.174}
\end{align}
（回忆对无界 $\phi$，有 $\alpha\geqslant0$：若 $\alpha<0$，则由 $\alpha\phi+\beta>0$ 得 $\phi\leqslant-\beta/\alpha$，与 $\phi$ 无界矛盾；而 $\alpha\geqslant-1$ 使 $\alpha+1>0$，故 $t_1^{\alpha+1}\to0$）。由\eqref{align:ch07-7.174}，序列 $\{u_{n_k}\}$ 对每个 $\tau\in(0,T)$ 含于 $L^\infty(\mathbb{R}\times(\tau,T))$ 的有界子集（因 $\phi$ 严格增且无界，$\phi^{-1}$ 在 $\phi$ 的值域上有定义，故 $u_{n_k}\leqslant\phi^{-1}((\alpha+1)\|u_0\|_{\mathscr{R}_f(\mathbb{R})}/\tau+\beta)$ 在 $\mathbb{R}\times(\tau,T)$ 上几乎处处成立），于是由\cref{theorem:ch07-7.3.1}，令 $n_k\to\infty$ 即得结论：极限密度满足同一 $L^\infty(\mathbb{R}\times(\tau,T))$ 界，且 $u_s$ 在 $\mathbb{R}\times(\tau,T)$ 上为零测度，故 $u=u_r$ 在该柱体上属于 $L^\infty$。
\end{proof}

% SUPPLEMENT: parent=proposition:ch07-7.3.5; kind=proof; 补出换元 $y=x-x_1$、$B_\varepsilon/\varepsilon$ 的极限（$Z^{-1}$ 在 $\lim_{y\to\infty}Z(y)$ 处的行为）、区间不相交性与递归终止的每一步，以及右端估计所用的\eqref{align:ch07-7.17}
\begin{proof}[\cref{proposition:ch07-7.3.5}的证明]
由\cref{theorem:ch07-7.3.3}的证明，对几乎所有 $t\in(0,T)$，不等式\eqref{align:ch07-7.172}对所有 $x_j\in\operatorname{supp}u_s(\cdot,t)$ 成立。设 $x_1\in\operatorname{supp}u_s(\cdot,t)$，置 $I_1\triangleq(x_1-\varepsilon,x_1+\varepsilon)$（$\varepsilon>0$）。因 $Z$ 可逆（$Z$ 把 $[0,\infty)$ 严格增地映到 $[0,\phi'(0))$，$\phi'(0)=\lim_{y\to\infty}Z(y)$），由\eqref{align:ch07-7.172}得，对 $x\in(x_1,x_1+\varepsilon)$ 且 $\lim_{y\to\infty}Z(y)-(x-x_1)/t\geqslant0$，
\begin{align*}
u_r(x,t)\geqslant Z^{-1}\!\left(\lim_{y\to\infty}Z(y)-\frac{x-x_1}{t}\right),
\end{align*}
故
\begin{align*}
\int_{I_1}u_r(x,t)\,\mathrm{d}x&\geqslant\int_{x_1}^{x_1+\varepsilon}Z^{-1}\!\left(\lim_{y\to\infty}Z(y)-\frac{x-x_1}{t}\right)\mathrm{d}x\\
&=\int_0^\varepsilon Z^{-1}\!\left(\lim_{y\to\infty}Z(y)-\frac{y}{t}\right)\mathrm{d}y=:B_\varepsilon,
\end{align*}
其中
\begin{align}
\lim_{\varepsilon\to0^+}\frac{B_\varepsilon}{\varepsilon}=Z^{-1}\!\left(\lim_{y\to\infty}Z(y)\right)=\infty, \label{align:ch07-7.175}
\end{align}
因为 $B_\varepsilon/\varepsilon$ 是 $y\mapsto Z^{-1}(\lim Z-y/t)$ 在 $(0,\varepsilon)$ 上的平均值，而该函数当 $y\to0^+$ 时趋于 $Z^{-1}(\lim Z)=\infty$（$Z^{-1}$ 在 $\lim_{y\to\infty}Z(y)=\sup Z$ 处无界）。

若 $\operatorname{supp}u_s(\cdot,t)\not\subseteq I_1$，取 $x_2\in\operatorname{supp}u_s(\cdot,t)\setminus I_1$ 并置 $I_2\triangleq(x_2-\varepsilon,x_2+\varepsilon)$。因 $x_2\notin I_1$ 即 $|x_1-x_2|\geqslant\varepsilon$，有 $(x_1,x_1+\varepsilon)\cap(x_2,x_2+\varepsilon)=\varnothing$，故 $u_r\geqslant0$ 给出
\begin{align*}
\int_{I_1\cup I_2}u_r(x,t)\,\mathrm{d}x\geqslant\int_{x_1}^{x_1+\varepsilon}u_r(x,t)\,\mathrm{d}x+\int_{x_2}^{x_2+\varepsilon}u_r(x,t)\,\mathrm{d}x\geqslant2B_\varepsilon.
\end{align*}
递归地，只要 $\operatorname{supp}u_s(\cdot,t)\not\subseteq\bigcup_{i=1}^{n-1}I_i$（$I_i\triangleq(x_i-\varepsilon,x_i+\varepsilon)$），就存在 $x_n\in\operatorname{supp}u_s(\cdot,t)\setminus\bigcup_{i=1}^{n-1}I_i$，于是半区间 $(x_i,x_i+\varepsilon)$（$i\leqslant n$）两两不交，且
\begin{align*}
nB_\varepsilon\leqslant\int_{\bigcup_{i=1}^nI_i}u_r(x,t)\,\mathrm{d}x\leqslant\|u_r(\cdot,t)\|_{L^1(\mathbb{R})}\leqslant\|u(\cdot,t)\|_{\mathscr{R}_f(\mathbb{R})}=\|u_0\|_{\mathscr{R}_f(\mathbb{R})},
\end{align*}
最后一步用\eqref{align:ch07-7.17}。故过程在某个 $n=n_\varepsilon$ 处终止，且 $n_\varepsilon B_\varepsilon\leqslant\|u_0\|_{\mathscr{R}_f(\mathbb{R})}$。于是 $\operatorname{supp}u_s(\cdot,t)\subseteq\bigcup_{i=1}^{n_\varepsilon}I_i$，从而
\begin{align*}
\lambda(\operatorname{supp}u_s(\cdot,t))\leqslant\lambda\!\left(\bigcup_{i=1}^{n_\varepsilon}I_i\right)\leqslant2n_\varepsilon\varepsilon\leqslant\frac{2\|u_0\|_{\mathscr{R}_f(\mathbb{R})}\,\varepsilon}{B_\varepsilon}.
\end{align*}
令 $\varepsilon\to0^+$ 并利用\eqref{align:ch07-7.175}，即得 $\lambda(\operatorname{supp}u_s(\cdot,t))=0$。结论得证。
\end{proof}

\end{document}
